% GENERATED FILE. Edit canonical lessons, then run npm run generate:books.
% canonical-source-hash: fnv1a64:498d72e8555f86d3
% canonical-lessons: UR-C421-kaghaz, UR-C421-kard, UR-C421-tokan, UR-C421-fis, UR-C421-karansi

\chapter{Papers, Cards and Fees}
\label{ch:ur-papers-cards-and-fees}
\hlchaptermodality{421}

\begin{chapteropening}
\textbf{By the end of this chapter:} \emph{I can deal with papers, cards, tokens and fees at a bank.}

Complete the last lesson of chapter 421: I can deal with papers, cards, tokens and fees at a bank.
\end{chapteropening}

\section[\texorpdfstring{\ur{کاغذ} (kāghaz)}{کاغذ (kāghaz)}]{\ur{کاغذ} (kāghaz) — paper; a document}

\label{lesson:UR-C421-kaghaz}



\begin{quote}
Before the new one: say the Urdu for to share out, then the Urdu for to cash (a cheque).
\end{quote}

\subsection*{You'll want to know: \ur{کاغذ}}
\textbf{\ur{کاغذ}} — \emph{kāghaz} — \textquotedblleft{}paper; a document\textquotedblright{}.

\begin{quote}
\ur{بینک} \ur{والے} \ur{نے} \ur{کچھ} \ur{کاغذ} \ur{مانگے۔}
\end{quote}

\begin{quote}
\emph{baink vāle ne kuchh kāghaz māṅge.}
\end{quote}

\begin{quote}
The man at the bank asked for some papers.
\end{quote}

Masculine.

\begin{cousinweb}
\textbf{\ur{کاغذ}} (\emph{kāghaz}) is usually traced to Persian \emph{kāghaz}, \textquotedblleft{}paper\textquotedblright{}; the word probably reached Persian from further east, along with paper itself.
\end{cousinweb}

\begin{grammarlens}[title={What you've built}]
One more word for getting around.
\end{grammarlens}

\subsection*{Your turn}
\begin{itemize}
\raggedright
  \item \emph{Say it:} \emph{kāghaz}
  \item \emph{Say it:} \emph{kāghaz}, once more
  \item \emph{Say it:} \emph{baink vāle ne kuchh kāghaz māṅge.} — The man at the bank asked for some papers.
  \item \emph{Recall:} say the Urdu for to share out, then the Urdu for to cash (a cheque), then say \emph{kāghaz} again
\end{itemize}

\begin{tcolorbox}[breakable,colback=teal!4,colframe=teal!35!black,title={Before you move on}]
What does \ur{کاغذ} mean? (\textquotedblleft{}Paper; a document\textquotedblright{}.) Say it once more.
\end{tcolorbox}
\section[\texorpdfstring{\ur{کارڈ} (kārḍ)}{کارڈ (kārḍ)}]{\ur{کارڈ} (kārḍ) — a card}

\label{lesson:UR-C421-kard}



\begin{quote}
Before the new one: say the Urdu for to cash (a cheque), then the Urdu for paper.
\end{quote}

\subsection*{You'll want to know: \ur{کارڈ}}
\textbf{\ur{کارڈ}} — \emph{kārḍ} — \textquotedblleft{}a card\textquotedblright{}.

\begin{quote}
\ur{میں} \ur{دکان} \ur{پر} \ur{کارڈ} \ur{سے} \ur{پیسے} \ur{دیتی} \ur{ہوں۔}
\end{quote}

\begin{quote}
\emph{maiṅ dukān par kārḍ se paise detī hūṅ.}
\end{quote}

\begin{quote}
At the shop I pay by card.
\end{quote}

Masculine.

\begin{cousinweb}
\textbf{\ur{کارڈ}} (\emph{kārḍ}) is borrowed from English \emph{card}.
\end{cousinweb}

\begin{grammarlens}[title={What you've built}]
One more word for getting around.
\end{grammarlens}

\subsection*{Your turn}
\begin{itemize}
\raggedright
  \item \emph{Say it:} \emph{kārḍ}
  \item \emph{Say it:} \emph{kārḍ}, once more
  \item \emph{Say it:} \emph{maiṅ dukān par kārḍ se paise detī hūṅ.} — At the shop I pay by card.
  \item \emph{Recall:} say the Urdu for to cash (a cheque), then the Urdu for paper, then say \emph{kārḍ} again
\end{itemize}

\begin{tcolorbox}[breakable,colback=teal!4,colframe=teal!35!black,title={Before you move on}]
What does \ur{کارڈ} mean? (\textquotedblleft{}A card\textquotedblright{}.) Say it once more.
\end{tcolorbox}
\section[\texorpdfstring{\ur{ٹوکن} (ṭokan)}{ٹوکن (ṭokan)}]{\ur{ٹوکن} (ṭokan) — a token}

\label{lesson:UR-C421-tokan}



\begin{quote}
Before the new one: say the Urdu for paper, then the Urdu for a card.
\end{quote}

\subsection*{You'll want to know: \ur{ٹوکن}}
\textbf{\ur{ٹوکن}} — \emph{ṭokan} — \textquotedblleft{}a token\textquotedblright{}.

\begin{quote}
\ur{بینک} \ur{میں} \ur{پہلے} \ur{ٹوکن} \ur{لیجیے۔}
\end{quote}

\begin{quote}
\emph{baink meṅ pehle ṭokan lījiye.}
\end{quote}

\begin{quote}
At the bank, please take a token first.
\end{quote}

Masculine.

\begin{cousinweb}
\textbf{\ur{ٹوکن}} (\emph{ṭokan}) is borrowed from English \emph{token}.
\end{cousinweb}

\begin{grammarlens}[title={What you've built}]
One more word for getting around.
\end{grammarlens}

\subsection*{Your turn}
\begin{itemize}
\raggedright
  \item \emph{Say it:} \emph{ṭokan}
  \item \emph{Say it:} \emph{ṭokan}, once more
  \item \emph{Say it:} \emph{baink meṅ pehle ṭokan lījiye.} — At the bank, please take a token first.
  \item \emph{Recall:} say the Urdu for paper, then the Urdu for a card, then say \emph{ṭokan} again
\end{itemize}

\begin{tcolorbox}[breakable,colback=teal!4,colframe=teal!35!black,title={Before you move on}]
What does \ur{ٹوکن} mean? (\textquotedblleft{}A token\textquotedblright{}.) Say it once more.
\end{tcolorbox}
\section[\texorpdfstring{\ur{فیس} (fīs)}{فیس (fīs)}]{\ur{فیس} (fīs) — a fee}

\label{lesson:UR-C421-fis}



\begin{quote}
Before the new one: say the Urdu for a card, then the Urdu for a token.
\end{quote}

\subsection*{You'll want to know: \ur{فیس}}
\textbf{\ur{فیس}} — \emph{fīs} — \textquotedblleft{}a fee\textquotedblright{}.

\begin{quote}
\ur{بینک} \ur{اس} \ur{کام} \ur{کی} \ur{فیس} \ur{لیتا} \ur{ہے۔}
\end{quote}

\begin{quote}
\emph{baink is kām kī fīs letā hai.}
\end{quote}

\begin{quote}
The bank charges a fee for this.
\end{quote}

Feminine.

\begin{cousinweb}
\textbf{\ur{فیس}} (\emph{fīs}) is borrowed from English \emph{fees}.
\end{cousinweb}

\begin{grammarlens}[title={What you've built}]
One more word for getting around.
\end{grammarlens}

\subsection*{Your turn}
\begin{itemize}
\raggedright
  \item \emph{Say it:} \emph{fīs}
  \item \emph{Say it:} \emph{fīs}, once more
  \item \emph{Say it:} \emph{baink is kām kī fīs letā hai.} — The bank charges a fee for this.
  \item \emph{Recall:} say the Urdu for a card, then the Urdu for a token, then say \emph{fīs} again
\end{itemize}

\begin{tcolorbox}[breakable,colback=teal!4,colframe=teal!35!black,title={Before you move on}]
What does \ur{فیس} mean? (\textquotedblleft{}A fee\textquotedblright{}.) Say it once more.
\end{tcolorbox}
\section[\texorpdfstring{\ur{کرنسی} (karansī)}{کرنسی (karansī)}]{\ur{کرنسی} (karansī) — a currency}

\label{lesson:UR-C421-karansi}



\begin{quote}
Before the new one: say the Urdu for a token, then the Urdu for a fee.
\end{quote}

\subsection*{You'll want to know: \ur{کرنسی}}
\textbf{\ur{کرنسی}} — \emph{karansī} — \textquotedblleft{}a currency\textquotedblright{}.

\begin{quote}
\ur{سرحد} \ur{پر} \ur{ہم} \ur{نے} \ur{کرنسی} \ur{بدلی۔}
\end{quote}

\begin{quote}
\emph{sarhad par ham ne karansī badlī.}
\end{quote}

\begin{quote}
We changed our money at the border.
\end{quote}

Feminine.

\begin{cousinweb}
\textbf{\ur{کرنسی}} (\emph{karansī}) is borrowed from English \emph{currency}.
\end{cousinweb}

\begin{grammarlens}[title={What you've built}]
One more word for getting around.
\end{grammarlens}

\subsection*{Your turn}
\begin{itemize}
\raggedright
  \item \emph{Say it:} \emph{karansī}
  \item \emph{Say it:} \emph{karansī}, once more
  \item \emph{Say it:} \emph{sarhad par ham ne karansī badlī.} — We changed our money at the border.
  \item \emph{Recall:} say the Urdu for a token, then the Urdu for a fee, then say \emph{karansī} again
\end{itemize}

\begin{tcolorbox}[breakable,colback=teal!4,colframe=teal!35!black,title={Before you move on}]
What does \ur{کرنسی} mean? (\textquotedblleft{}A currency\textquotedblright{}.) Say it once more.
\end{tcolorbox}
